\documentclass[12pt,a4paper]{report}

% ===== Packages =====
\usepackage[utf8]{inputenc}
\usepackage[T1]{fontenc}
\usepackage[french]{babel}
\usepackage{csquotes} % recommandé avec biblatex
\usepackage{graphicx}
\usepackage{titlesec}
\usepackage{fancyhdr}
\usepackage{setspace}
\usepackage{float}
\usepackage{caption}
\usepackage{appendix}
\usepackage{hyperref}
\usepackage{tocloft}
\usepackage[table]{xcolor}
\usepackage{array}
\usepackage{etoolbox} % utile pour certains ajustements
\usepackage{placeins}
% Bibliographie
\usepackage{enumitem}
\usepackage[backend=biber,style=numeric,sorting=none]{biblatex}
\addbibresource{webliography.bib}
 
\usepackage{subcaption}
% Glossaire / acronymes
\usepackage[acronym]{glossaries}
\makeglossaries

% ===== Marges =====
\usepackage[a4paper,
            left=3cm,
            right=2.5cm,
            top=2.5cm,
            bottom=2.5cm,
            headheight=14pt,
            headsep=0.8cm,
            footskip=1.2cm]{geometry}

% ===== Interligne =====
\onehalfspacing

% ===== En-têtes / pieds =====
\pagestyle{fancy}
\fancyhf{}
\fancyhead[L]{\small\scshape\leftmark} % plus petit et petites majuscules
\fancyhead[R]{\thepage}
\renewcommand{\headrulewidth}{0.4pt}
\renewcommand{\footrulewidth}{0pt}

% ===== Titres =====
% Chapitre
\titleformat{\chapter}[display]
  {\normalfont\huge\bfseries\centering}
  {\chaptertitlename\ \thechapter}
  {20pt}
  {\Huge}
\titlespacing*{\chapter}{0pt}{-10pt}{30pt} % espacement avant/après

% Section
\titleformat{\section}
  {\normalfont\Large\bfseries}
  {\thesection}{1em}{}
\titlespacing*{\section}{0pt}{3ex plus .5ex minus .2ex}{2ex plus .2ex}

% Sous-section
\titleformat{\subsection}
  {\normalfont\large\bfseries}
  {\thesubsection}{1em}{}
\titlespacing*{\subsection}{0pt}{2.5ex plus .5ex minus .2ex}{1.5ex plus .2ex}

% ===== Légendes =====
\captionsetup{
    font=small,
    labelfont=bf,
    labelsep=period
}
\renewcommand{\figurename}{Figure}
\renewcommand{\tablename}{Tableau}

% ===== Table des matières =====
\setcounter{tocdepth}{3}
\setcounter{secnumdepth}{3}
\renewcommand{\cftchapleader}{\cftdotfill{\cftdotsep}}
\renewcommand{\cftsecleader}{\cftdotfill{\cftdotsep}}

% ===== Liens =====
\hypersetup{
    colorlinks=true,
    linkcolor=black,
    filecolor=black,      
    urlcolor=blue,
    citecolor=black,
    bookmarks=true,
    bookmarksopen=true,
    pdftitle={Rapport PFA - AutoStory - Générateur intelligent de contenus publicitaires pour le secteur automobile},
    pdfauthor={Youssef MOUSTAID et Fayz OUSSAMA}
}

% ===== Titre =====
\title{
  Rapport de Projet de Fin d'Année \\[0.5em]
  \Large AutoStory - Générateur intelligent de contenus publicitaires pour le secteur automobile
}
\author{
  Youssef MOUSTAID \\[0.2em]
  Fayz OUSSAMA \\[0.5em]
  \small Filière : Méthodes Informatiques Appliquées à la Gestion des Entreprises - MIAGE  \\
  \small 9ème Semestre \\[0.5em]
  \small Encadré par : M. RAOUYANE  \\
  \small Année Universitaire 2025-2026
}
\date{30 décembre 2025}

% ===== Document =====
\begin{document}

% Page de garde
\begin{titlepage}
    \centering
    \includegraphics[width=0.35\textwidth]{images/logos/emsi_logo.png}\par\vspace{1cm}
    {\scshape\Large EMSI \par}
    \vspace{0.5cm}
    {\scshape\large Filière : Méthodes Informatiques Appliquées à la Gestion des Entreprises - MIAGE\par}
    {\scshape\large 9ème Semestre\par}
    \vspace{1.5cm}
    {\huge\bfseries Rapport de Projet de Fin d'Année\par}
    \vspace{0.5cm}
    
    % Logo AutoStory
    \includegraphics[width=0.1\textwidth]{images/autostory-logo.png}\par
    \vspace{0.5cm}
    
    {\Large\bfseries AutoStory\par}
    \vspace{0.3cm}
    {\large Générateur intelligent de contenus publicitaires\par}
    {\large pour le secteur automobile\par}
    \vspace{1.5cm}
    {\Large Réalisé par\par}
    \vspace{0.5cm}
    {\large 
        \begin{tabular}{c}
            Youssef MOUSTAID \\
            Fayz OUSSAMA
        \end{tabular}
    \par}
    \vfill
    Encadré par\par
    M. RAOUYANE
    \vfill
    {\large Année Universitaire 2025-2026\par}
    {\large Décembre 2025\par}
\end{titlepage}

% Page de validation
\clearpage








% dedicases
\chapter*{Dédicaces}
\addcontentsline{toc}{chapter}{Dédicaces}

Je tiens à dédier ce travail à toutes les personnes qui ont marqué mon parcours et qui m’ont permis de mener à bien ce \textbf{Projet de Fin d’Année}.  

\vspace{1em}

\textbf{À mes chers parents},  
pour leur amour incommensurable, leur patience et leurs sacrifices. Leur soutien moral, matériel et affectif a été la clé de ma réussite. Je leur dois ce que je suis aujourd’hui et leur exprime ma gratitude éternelle.  

\vspace{1em}

\textbf{À ma famille},  
frères, sœurs et proches, pour leurs encouragements constants, leur compréhension et leurs prières qui m’ont donné la force de persévérer même dans les moments difficiles.  

\vspace{1em}

\textbf{À mes enseignants et encadrants},  
qui m’ont transmis le savoir, guidé dans ma formation et accompagné tout au long de ce projet. Leur disponibilité et leurs conseils avisés ont grandement contribué à l’aboutissement de ce travail.  

\vspace{1em}

\textbf{À mes camarades de promotion et amis},  
avec qui j’ai partagé des moments de travail, de réflexion et d’amitié. Leur collaboration et leur esprit d’équipe ont rendu ce parcours plus agréable et enrichissant.  

\vspace{1em}

\textbf{Enfin, à toutes les personnes}, proches ou lointaines, qui ont cru en moi et m’ont soutenu dans cette aventure académique et personnelle. Ce travail est aussi le leur.  

% Remerciements
\chapter*{Remerciements}
\addcontentsline{toc}{chapter}{Remerciements}
Nous tenons à exprimer notre profonde gratitude à toutes les personnes qui ont contribué de près ou de loin à la réalisation de ce \textbf{Projet de Fin d'Année}.  
\vspace{1em}

Tout d'abord, nous remercions l'\textbf{École Marocaine des Sciences de l'Ingénieur (EMSI)} pour la formation de qualité dont nous avons bénéficié et pour nous avoir offert l'opportunité de travailler sur un projet innovant dans le domaine de l'intelligence artificielle appliquée au marketing automobile.
\vspace{1em}

Nous souhaitons également remercier notre encadrant académique  \textbf{M. Raouyane}  pour son suivi rigoureux, ses conseils avisés et ses recommandations constructives, qui nous ont permis d'améliorer la qualité de ce travail et de mener à bien ce projet.  
\vspace{1em}

Nos remerciements vont également à tous les enseignants de la filière \textbf{MIAGE} qui nous ont transmis les connaissances théoriques et pratiques nécessaires à la réalisation de ce projet.
\vspace{1em}

Enfin, nous adressons un grand merci à nos familles et nos amis pour leur soutien moral et leur encouragement constant durant cette période exigeante.

% Résumé
\chapter*{Résumé}
\addcontentsline{toc}{chapter}{Résumé}
Ce rapport présente le \textbf{Projet de Fin d'Année (PFA)} intitulé \textit{AutoStory}, réalisé à l'\textbf{École Marocaine des Sciences de l'Ingénieur (EMSI)} dans le cadre de la filière \textbf{MIAGE}. L'objectif principal de ce projet est le développement d'un \textbf{système automatisé de génération de contenus publicitaires multimédia} pour le secteur automobile, exploitant les capacités des \textbf{Large Language Models (LLMs)} et des technologies web modernes.
\vspace{1em}

La solution développée se compose d'un \textbf{backend Node.js/Express} utilisant \textbf{Google Gemini AI 2.5 Flash} pour la génération de scripts narratifs techniques, \textbf{Puppeteer} pour la composition vidéo automatisée, et \textbf{MongoDB} pour le stockage des données. Le \textbf{frontend React/TypeScript} offre une interface utilisateur moderne et intuitive avec des visualisations 3D interactives via \textbf{Three.js}. Le système permet de générer des contenus publicitaires personnalisés (texte, vidéo, PDF, 3D) adaptés au style du véhicule.
\vspace{1em}

Ce projet académique a permis d'acquérir une expérience enrichissante en matière d'\textbf{intelligence artificielle générative}, de \textbf{développement full-stack}, d'\textbf{architecture logicielle modulaire}, et de traitement automatisé du langage naturel, tout en développant une solution innovante pour le marketing automobile.
\vspace{1em}

\bigskip
\textbf{Mots-clés :} Intelligence Artificielle, LLM, Gemini AI, Node.js, JavaScript, React, Génération automatique de contenu, Marketing automobile, Three.js, MongoDB.

% TOC et listes
\clearpage
\tableofcontents

\clearpage
\listoffigures
\addcontentsline{toc}{chapter}{Liste des figures}

\clearpage
\listoftables
\addcontentsline{toc}{chapter}{Liste des tableaux}

\clearpage
\chapter*{Liste des sigles et abréviations}
\addcontentsline{toc}{chapter}{Liste des sigles et abréviations}
\printglossary[type=\acronymtype,title=Liste des sigles et abréviations]

% --- Contenu ---

\chapter{Introduction}
\section*{Introduction Générale}

Ce rapport de \textbf{Projet de Fin d'Année (PFA)} présente le développement du projet \textbf{AutoStory}, réalisé dans le cadre de la formation en \textbf{Ingénierie Logicielle et Intelligence Artificielle} à l'\textbf{École Marocaine des Sciences de l'Ingénieur (EMSI)}.

\subsection*{Contexte}

Le marketing digital et le secteur automobile connaissent actuellement une transformation profonde portée par l'essor des technologies d'\textbf{intelligence artificielle générative}. Les entreprises du secteur automobile ont un besoin croissant de \textbf{contenus publicitaires personnalisés et produits rapidement} pour promouvoir leurs véhicules sur les plateformes digitales. Cependant, la production vidéo traditionnelle reste \textbf{coûteuse, chronophage et nécessite des compétences spécialisées} (scénaristes, voix-off professionnelles, monteurs vidéo).

L'émergence des \textbf{Large Language Models (LLMs)} tels que GPT-4, Claude et Google Gemini ouvre de nouvelles perspectives pour l'automatisation de la création de contenu narratif. Parallèlement, les technologies de \textbf{synthèse vocale (Text-to-Speech)} et de \textbf{composition vidéo programmatique} permettent aujourd'hui de produire des supports publicitaires de qualité sans intervention humaine extensive.

\subsection*{Problématique}

Face à ce contexte, plusieurs questions se posent :

\begin{itemize}
    \item Comment exploiter les capacités des LLMs pour générer des scripts publicitaires techniques et cohérents pour des véhicules automobiles ?
    \item Comment automatiser l'ensemble de la chaîne de production vidéo, de la génération du texte à la composition finale ?
    \item Comment garantir la qualité narrative et la cohérence des contenus générés automatiquement ?
    \item Quelle architecture logicielle adopter pour assurer modularité, scalabilité et maintenabilité du système ?
\end{itemize}

\subsection*{Objectifs du projet}

Le projet \textbf{AutoStory} vise à répondre à ces problématiques en développant un système complet de génération automatisée de vidéos publicitaires pour le secteur automobile. Les objectifs spécifiques sont les suivants :

\begin{enumerate}
    \item Mettre en place un système de génération de scripts narratifs techniques utilisant \textbf{Google Gemini 2.5 Flash}
    \item Développer un module de synthèse vocale en français avec \textbf{gTTS (Google Text-to-Speech)}
    \item Implémenter un système de composition vidéo automatisé avec \textbf{MoviePy}
    \item Créer un module de classification automatique des styles de véhicules (sportif, luxe, familial, écologique)
    \item Concevoir une interface utilisateur moderne avec \textbf{React et TypeScript}
    \item Adopter une architecture logicielle \textbf{modulaire et scalable} facilitant l'évolution future du système
\end{enumerate}

\subsection*{Démarche et structure du rapport}

Ce travail s'inscrit dans une démarche pédagogique rigoureuse, combinant apprentissage théorique et application pratique. Il nous permet de mobiliser nos connaissances en \textbf{intelligence artificielle}, \textbf{développement logiciel}, \textbf{traitement du langage naturel} et \textbf{architecture applicative}.

Le présent rapport est structuré de la manière suivante :

\begin{itemize}
    \item Le \textbf{Chapitre 1} présente le cahier des charges du projet AutoStory
    \item Le \textbf{Chapitre 2} expose le contexte général, l'état de l'art des technologies IA génératives et le positionnement d'AutoStory
    \item Le \textbf{Chapitre 3} détaille l'analyse et la spécification des besoins fonctionnels et non-fonctionnels
    \item Le \textbf{Chapitre 4} décrit la conception de l'architecture et les choix techniques
    \item Le \textbf{Chapitre 5} présente la réalisation et la mise en œuvre du système
    \item Le \textbf{Chapitre 6} conclut par un bilan du projet et présente les perspectives d'évolution
\end{itemize}

\chapter{Contexte Général et Présentation du Projet MobilLoca}
\label{chap:contexte}

\section{Introduction}
\label{sec:intro-contexte}
Dans ce chapitre, nous présenterons le cadre général de notre projet de fin d'études réalisé au sein de l'entreprise \textbf{Marketing Confort}, un acteur dynamique spécialisé dans le marketing digital et le développement de solutions technologiques innovantes. Nous débuterons par une présentation détaillée de l'organisme d'accueil, de ses activités, de sa vision et de ses compétences techniques. Ensuite, nous introduirons le projet \textbf{MobilLoca}, en explicitant son contexte, sa problématique, ses objectifs stratégiques et fonctionnels, ainsi que son plan de réalisation. Ce chapitre vise à poser les bases contextuelles et organisationnelles nécessaires à la compréhension des enjeux et de la portée du projet.

\section{L'organisme d'accueil}
\label{sec:organisme-accueil}

\subsection{Présentation de l'organisation}
\label{subsec:presentation-organisation}

\textbf{Marketing Confort} est une agence de marketing digital et de développement web fondée en 2013 et basée à Fès, avec une présence également en France. Elle accompagne ses clients dans la conception et la réalisation de projets digitaux sur mesure, allant de la stratégie marketing à la production audiovisuelle, en passant par le développement d'applications web et mobiles. L'entreprise se distingue par son approche personnalisée et son souci constant d'innovation, visant à offrir à ses clients un avantage concurrentiel durable grâce à des solutions digitales performantes et adaptées.

\begin{figure}[H]
    \centering
    \begin{minipage}{0.3\textwidth}
        \centering
        \includegraphics[width=\linewidth]{images/logos/marketing comfort.png}
        \caption{Logo de Marketing Confort}
        \label{fig:logo_marketing_confort}
    \end{minipage}
    \hfill
    \begin{minipage}{0.65\textwidth}
        \subsection{Mission et vision}
        \label{subsec:mission-vision}
        La mission de \textbf{Marketing Confort} est de rendre le marketing digital accessible et performant pour tous types d'entreprises, quelle que soit leur taille ou leur secteur d'activité. 
        Sa vision est de devenir un leader mondial dans le domaine du marketing numérique et de la production audiovisuelle, en établissant de nouveaux standards d'excellence et d'innovation.
    \end{minipage}
\end{figure}

\subsection{Services proposés}
\label{subsec:services-proposes}

L'agence propose une gamme étendue de services, incluant :

\begin{itemize}
    \item \textbf{Développement Web \& Mobile} : Conception et développement d'applications web et mobiles sur mesure, optimisées pour l'expérience utilisateur et le référencement.
    \item \textbf{Marketing Digital} : Stratégies de communication, campagnes publicitaires (Google Ads, Facebook Ads), gestion des réseaux sociaux et optimisation \gls{seo}.
    \item \textbf{Production Audiovisuelle} : Création de contenus vidéo promotionnels, spots publicitaires et supports visuels pour renforcer l'image de marque.
\end{itemize}

\begin{table}[H]
\centering
\renewcommand{\arraystretch}{1.5}
\caption{Fiche d'identité de Marketing Confort}
\label{tab:fiche_identite_mc}
\begin{tabular}{|>{\bfseries}p{0.3\textwidth}|p{0.6\textwidth}|}
\hline
\rowcolor{blue!25} \textbf{Information} & \textbf{Détail} \\
\hline
Nom de la société & Marketing Confort \\
\hline
ICE & 002720451000006 \\
\hline
Forme juridique & Société à Responsabilité Limitée \\
\hline
Date de création & 2013 \\
\hline
Siège Social & N°22 Lotissement Idrissi 3ème étage Atlas, Fès, Maroc \\
\hline
Coordonnées & +212 612 14 06 66 / +33760-297926 \\
 & contact@marketingconfort.com \\
\hline
Secteur d'activité & Services d'information commerciale \\
\hline
Nombre d'employés & 11 - 50 Employés \\
\hline
Capital & 100 000 DH \\
\hline
\end{tabular}
\end{table}
\cite{charika2025marketing}


\subsection{Partenaires stratégiques}
\label{subsec:partenaires-strategiques}

Marketing Confort collabore avec plusieurs partenaires prestigieux, tels que :

\begin{itemize}
    \item AMDL (Agence Marocaine de Développement de la Logistique)
    \item Eau Thermale Avène
    \item TotalEnergies
    \item Burger King
    \item Bridges to the Future
    \item Minéraux Maroc
    \item Le Phare Pressing
    \item iCi (La Réussite)
    \item Isyane (Senteurs et Saveurs)
\end{itemize}

Ces collaborations enrichissent son expertise et lui permettent d'intervenir dans des secteurs variés.

\subsection{Structure organisationnelle}
\label{subsec:structure-organisationnelle}

L'entreprise est dirigée par un CEO et un CTO, et s'appuie sur des équipes spécialisées en développement IT, marketing, production audiovisuelle, ainsi que sur un réseau de freelances et de partenaires.

\begin{figure}[H]
    \centering
    \includegraphics[width=0.7\textwidth]{images/chapitre 2 - contexte/organigrame.png}
    \caption{Organigramme de Marketing Confort}
    \label{fig:organigramme_mc}
\end{figure}

\subsection{Compétences techniques}
\label{subsec:competences-techniques}

Marketing Confort dispose de compétences techniques diversifiées, couvrant :

\begin{itemize}
    \item Développement Web \& Mobile
    \item Cloud Computing
    \item Marketing Digital \& Production Audiovisuelle
    \item Coaching Agile \& DevOps
    \item Intelligence Artificielle
\end{itemize}

\begin{figure}[H]
    \centering
    \includegraphics[width=0.8\textwidth]{images/chapitre 2 - contexte/competences.png}
    \caption{Compétences de l'entreprise}
    \label{fig:competences_mc}
\end{figure}

\section{Présentation du projet MobilLoca}
\label{sec:presentation-projet}

\subsection{Contexte du projet}
\label{subsec:contexte-projet}

Le projet \textbf{MobilLoca} a été initié à la demande d'\textbf{Adanev}, une entreprise française spécialisée dans le transport de personnes à mobilité réduite. Bien qu'Adanev dispose déjà d'une application de transport avec chauffeur (\textit{Adanev Cab}), elle souhaitait étendre son offre en proposant un service de location de véhicules sans chauffeur, avec une gestion centralisée et numérique de l'ensemble du processus.

\subsection{Problématique}
\label{subsec:problematique}

La gestion manuelle des réservations, de la flotte, des documents administratifs et des paiements entraînait plusieurs limitations :

\begin{itemize}
    \item Manque de visibilité sur les performances globales
    \item Processus longs et sujets aux erreurs
    \item Expérience client peu fluide et manque d'autonomie
    \item Difficulté à suivre les maintenances et les documents expirés
\end{itemize}

\subsection{Objectifs du projet}
\label{subsec:objectifs-projet}

\textbf{MobilLoca} vise à :

\begin{itemize}
    \item \textbf{Automatiser} la gestion des réservations, contrats, paiements et documents
    \item \textbf{Centraliser} la gestion de la flotte et des agences
    \item \textbf{Offrir une expérience client intuitive} via une application mobile
    \item \textbf{Fournir des outils d'aide à la décision} via des tableaux de bord et rapports analytiques
    \item \textbf{Garantir une solution scalable et sécurisée}
\end{itemize}

\subsection{Solution proposée}
\label{subsec:solution-proposee}

La solution retenue est une plateforme digitale complète comprenant :

\begin{itemize}
    \item Un \textbf{Back-Office Web} pour les agents : gestion des véhicules, réservations, documents, paiements, statistiques.
    \item Une \textbf{Application Mobile} pour les clients : consultation du catalogue, réservation en ligne, paiement sécurisé, gestion des documents et historique.
\end{itemize}

\section{Planification du projet}
\label{sec:planification-projet}

\subsection{Méthodologie de gestion de projet}
\label{subsec:methodologie-gestion}

Une approche \textbf{Agile Scrum} a été adoptée pour garantir flexibilité, réactivité et livraison incrémentale. Le projet a été découpé en sprints de 2 à 3 semaines, avec des revues régulières et des ajustements continus.

\subsection{Phases principales}
\label{subsec:phases-principales}

\begin{enumerate}
    \item \textbf{Spécifications et conception} (Février 03 - Mai 01 2025)
    \item \textbf{Développement Front-End Web} (Mars 03 - Mai 16 2025)
    \item \textbf{Développement Front-End Mobile} (Février 16 - Mai 16 2025)
    \item \textbf{Structuration} (Avril 17 - Mai 17 2025)
    \item \textbf{Initiation des micro-services Back-End} (Mai 19 - Juin 02 2025)
    \item \textbf{Développement des micro-services Back-End} (Mai 19 - Septembre 02 2025)
    \item \textbf{Intégration des micro-service dans les applications} (Juilled 03 - Septembre 02 2025)
\end{enumerate}

\begin{figure}[H]
    \centering
    \includegraphics[width=\textwidth]{images/chapitre 2 - contexte/gant.png}
    \caption{Diagramme de Gantt du projet MobilLoca}
    \label{fig:gantt_mobilloca}
\end{figure}
\begin{figure}[H]
    \centering
    \includegraphics[width=\textwidth]{images/chapitre 2 - contexte/mobiloca_timeline.png}
    \caption{Diagramme de Gantt du projet MobilLoca - Derniers phases }
    \label{fig:gantt_mobilloca}
\end{figure}

\subsection{Outils de collaboration}
\subsubsection{Les outils :}
\label{subsec:outils-collaboration}

Pour garantir une communication efficace et une gestion de projet optimale, nous avons utilisé plusieurs outils collaboratifs adaptés à nos besoins. Chaque outil joue un rôle spécifique dans le processus de développement, permettant à l'équipe de rester synchronisée et productive.

\vspace{0.5em} % petit espace avant la liste

% Jira
\noindent
\begin{minipage}[H]{0.1\textwidth}
    \includegraphics[width=\linewidth]{images/logos/jira.png}
\end{minipage}%
\hfill
\begin{minipage}[H]{0.85\textwidth}
    \textbf{Jira} : Outil de gestion des tâches et des sprints, Jira nous a permis de planifier, suivre et prioriser les activités de développement selon la méthodologie Agile Scrum. Chaque tâche, user story ou bug était clairement assigné, ce qui garantissait transparence et traçabilité tout au long du projet.
\end{minipage}

\vspace{0.5em} % espace entre les outils

% GitLab
\noindent
\begin{minipage}[H]{0.1\textwidth}
    \includegraphics[width=\linewidth]{images/logos/gitlab.png}
\end{minipage}%
\hfill
\begin{minipage}[H]{0.85\textwidth}
    \textbf{GitLab} : Nous avons utilisé GitLab pour le versioning et l'intégration continue. Chaque membre de l'équipe pouvait gérer son code via des branches, des commits et des merge requests. Les pipelines CI/CD automatisés ont permis de tester et déployer les applications sans erreurs, améliorant la qualité globale du projet.
\end{minipage}

\vspace{0.5em}

% Discord
\noindent
\begin{minipage}[H]{0.1\textwidth}
    \includegraphics[width=\linewidth]{images/logos/discord.png}
\end{minipage}%
\hfill
\begin{minipage}[H]{0.85\textwidth}
    \textbf{Discord} : Pour la communication en temps réel, nous avons utilisé Discord, ce qui a facilité les discussions instantanées, les réunions d'équipe virtuelles et le partage de fichiers. L'outil a contribué à maintenir un flux de communication constant entre tous les membres.
\end{minipage}

\vspace{0.5em}

% Titan-Mail
\noindent
\begin{minipage}[H]{0.1\textwidth}
    \includegraphics[width=\linewidth]{images/logos/titan.png}
\end{minipage}%
\hfill
\begin{minipage}[H]{0.85\textwidth}
    \textbf{Titan-Mail} : Titan-Mail a été utilisé pour la messagerie professionnelle, garantissant la gestion formelle des communications, le suivi des échanges et la réception des notifications importantes liées au projet.
\end{minipage}


\subsubsection{Extraits et exemples :}

% === Jira examples ===
\noindent
\textbf{Jira :} Exemples de gestion des sprints et des tâches.

\begin{figure}[H]
    \centering
    \includegraphics[width=0.9\textwidth]{images/chapitre 2 - contexte/board.png}
    \caption{Exemple d’un sprint board dans Jira}
    \label{fig:jira-sprint-board}
\end{figure}

\vspace{1em}

% === GitLab examples ===
\noindent
\textbf{GitLab :} Exemples de gestion du code, merge requests et pipelines CI/CD.

\begin{figure}[H]
    \centering
    \includegraphics[width=0.85\textwidth]{images/chapitre 2 - contexte/merge-request.png}
    \caption{Exemple de merge request dans GitLab}
    \label{fig:gitlab-mr}
\end{figure}

\begin{figure}[H]
    \centering
    \includegraphics[width=0.85\textwidth]{images/chapitre 2 - contexte/git-lab-modules.png}
    \caption{Modules de projet et structure dans GitLab}
    \label{fig:gitlab-modules}
\end{figure}

\vspace{1em}

% === Discord examples ===
\noindent
\textbf{Discord :} Communication en temps réel et suivi des discussions de projet.

\begin{figure}[H]
    \centering
    \includegraphics[width=0.85\textwidth]{images/chapitre 2 - contexte/asking-devops-engineer-to-validate-mr.png}
    \caption{Exemple de demandes et discussions dans Discord}
    \label{fig:discord-requests}
\end{figure}

\vspace{1em}


\section{Conclusion}
\label{sec:conclusion-contexte}

Ce chapitre a permis de situer le projet \textbf{MobilLoca} dans son contexte organisationnel et stratégique. Nous avons présenté l'entreprise d'accueil, ses services, sa vision, ainsi que les enjeux et objectifs du projet. La méthodologie Agile et les outils de gestion mis en place assurent un cadre de travail structuré et réactif. Le chapitre suivant détaillera l'analyse des besoins et la conception technique de l'application.
\chapter{Analyse et Spécification des Besoins}

\section{Introduction}

Dans ce chapitre, nous présenterons l'analyse détaillée des besoins du projet AutoStory. Cette phase constitue une étape fondamentale dans le développement du système, car elle permet d'identifier et de formaliser les exigences fonctionnelles et non-fonctionnelles. Nous établirons également les spécifications techniques qui guideront la conception et l'implémentation de la solution.

\section{Besoins fonctionnels}
\label{sec:besoins-fonctionnels}

Les besoins fonctionnels décrivent les fonctionnalités que le système AutoStory doit offrir aux utilisateurs.

\subsection{Génération de scripts narratifs}
\label{subsec:generation-scripts}

\textbf{BF1 : Génération automatique de scripts publicitaires}
\begin{itemize}
    \item Le système doit générer des scripts narratifs techniques et cohérents pour des véhicules automobiles
    \item Le script doit inclure des informations sur les caractéristiques techniques, le design et les innovations
    \item La longueur du script doit être adaptée à une vidéo de 30-60 secondes
    \item Le ton et le style doivent s'adapter au type de véhicule (sportif, luxe, familial, écologique)
\end{itemize}

\textbf{BF2 : Utilisation d'un Large Language Model}
\begin{itemize}
    \item Le système doit intégrer Google Gemini 2.5 Flash pour la génération de textes
    \item Les prompts doivent être optimisés pour obtenir des résultats de qualité professionnelle
    \item Le système doit gérer les erreurs de communication avec l'API Gemini
\end{itemize}

\subsection{Classification des véhicules}
\label{subsec:classification-vehicules}

\textbf{BF3 : Détection automatique du style de véhicule}
\begin{itemize}
    \item Le système doit classifier automatiquement les véhicules en catégories (sportif, luxe, familial, écologique)
    \item La classification doit se baser sur les caractéristiques techniques et visuelles du véhicule
    \item Le système doit adapter le ton narratif en fonction de la catégorie détectée
\end{itemize}

\subsection{Synthèse vocale}
\label{subsec:synthese-vocale}

\textbf{BF4 : Conversion texte vers audio}
\begin{itemize}
    \item Le système doit convertir le script généré en audio avec une voix française naturelle
    \item La qualité audio doit être professionnelle (format MP3, bitrate adapté)
    \item La synthèse vocale doit gérer correctement la prononciation des termes techniques automobiles
\end{itemize}

\subsection{Composition vidéo}
\label{subsec:composition-video}

\textbf{BF5 : Génération automatique de vidéos}
\begin{itemize}
    \item Le système doit composer automatiquement une vidéo à partir d'images du véhicule
    \item La vidéo doit synchroniser les images avec la narration audio
    \item Le système doit ajouter des transitions fluides entre les images
    \item La vidéo finale doit être exportée en format MP4 haute définition (1080p minimum)
\end{itemize}

\textbf{BF6 : Personnalisation visuelle}
\begin{itemize}
    \item Le système doit permettre l'ajout de textes overlay (nom du véhicule, slogan)
    \item Les effets visuels doivent correspondre au style du véhicule
    \item Le système doit gérer différents formats d'images en entrée
\end{itemize}

\subsection{Interface utilisateur}
\label{subsec:interface-utilisateur}

\textbf{BF7 : Interface web intuitive}
\begin{itemize}
    \item L'utilisateur doit pouvoir uploader des images du véhicule
    \item L'utilisateur doit pouvoir saisir les caractéristiques techniques du véhicule
    \item Le système doit afficher l'état d'avancement de la génération (script, audio, vidéo)
    \item L'utilisateur doit pouvoir prévisualiser et télécharger la vidéo générée
\end{itemize}

\textbf{BF8 : Gestion de multiples générations}
\begin{itemize}
    \item L'utilisateur doit pouvoir consulter l'historique des vidéos générées
    \item Le système doit permettre de régénérer une vidéo avec des paramètres modifiés
    \item L'utilisateur doit pouvoir supprimer des vidéos de son historique
\end{itemize}

\section{Besoins non-fonctionnels}
\label{sec:besoins-non-fonctionnels}

Les besoins non-fonctionnels définissent les contraintes et les critères de qualité du système.

\subsection{Performance}
\label{subsec:performance}

\textbf{BNF1 : Temps de génération}
\begin{itemize}
    \item Le temps total de génération d'une vidéo ne doit pas excéder 5 minutes
    \item La génération du script doit prendre moins de 30 secondes
    \item La synthèse vocale doit prendre moins de 20 secondes
    \item La composition vidéo doit prendre moins de 3 minutes
\end{itemize}

\textbf{BNF2 : Optimisation des ressources}
\begin{itemize}
    \item Le système doit optimiser l'utilisation de la mémoire lors du traitement vidéo
    \item Les fichiers temporaires doivent être nettoyés après la génération
\end{itemize}

\subsection{Fiabilité}
\label{subsec:fiabilite}

\textbf{BNF3 : Gestion des erreurs}
\begin{itemize}
    \item Le système doit gérer gracieusement les échecs d'API (Gemini, gTTS)
    \item Les erreurs doivent être loggées avec des informations détaillées
    \item L'utilisateur doit recevoir des messages d'erreur clairs et actionnables
\end{itemize}

\textbf{BNF4 : Qualité des contenus}
\begin{itemize}
    \item Les scripts générés doivent être cohérents et sans erreurs grammaticales
    \item La qualité audio doit être professionnelle sans distorsions
    \item La qualité vidéo doit être HD (1920x1080 minimum)
\end{itemize}

\subsection{Maintenabilité}
\label{subsec:maintenabilite}

\textbf{BNF5 : Architecture modulaire}
\begin{itemize}
    \item Le code doit être organisé en modules indépendants (génération, synthèse, composition)
    \item Chaque module doit avoir une responsabilité unique et bien définie
    \item Le système doit permettre le remplacement facile d'un composant (ex: changer de LLM)
\end{itemize}

\textbf{BNF6 : Documentation}
\begin{itemize}
    \item Le code doit être commenté et documenté (docstrings Python)
    \item Un guide d'utilisation doit être fourni pour les utilisateurs
    \item Une documentation technique doit être disponible pour les développeurs
\end{itemize}

\subsection{Sécurité}
\label{subsec:securite}

\textbf{BNF7 : Protection des données}
\begin{itemize}
    \item Les clés API doivent être stockées de manière sécurisée (variables d'environnement)
    \item Les fichiers uploadés doivent être validés (type, taille, contenu)
    \item Les données utilisateur ne doivent pas être partagées avec des tiers
\end{itemize}

\subsection{Scalabilité}
\label{subsec:scalabilite}

\textbf{BNF8 : Évolutivité}
\begin{itemize}
    \item L'architecture doit permettre l'ajout de nouveaux styles de véhicules
    \item Le système doit supporter l'ajout de nouvelles langues pour la synthèse vocale
    \item Le système doit être prêt pour une migration vers le cloud si nécessaire
\end{itemize}

\subsection{Utilisabilité}
\label{subsec:utilisabilite}

\textbf{BNF9 : Expérience utilisateur}
\begin{itemize}
    \item L'interface doit être intuitive et ne nécessiter aucune formation
    \item Le système doit fournir des feedbacks visuels clairs sur l'état du traitement
    \item L'interface doit être responsive et accessible sur différents appareils
\end{itemize}

\section{Contraintes techniques}
\label{sec:contraintes-techniques}

\subsection{Contraintes technologiques}
\begin{itemize}
    \item \textbf{Backend :} Python 3.10+, FastAPI
    \item \textbf{Frontend :} React 18+, TypeScript 5+
    \item \textbf{LLM :} Google Gemini 2.5 Flash
    \item \textbf{TTS :} gTTS (Google Text-to-Speech)
    \item \textbf{Composition vidéo :} MoviePy, Pillow
    \item \textbf{Format vidéo :} MP4 (H.264)
\end{itemize}

\subsection{Contraintes d'environnement}
\begin{itemize}
    \item Le système doit fonctionner sur Linux et Windows
    \item Les dépendances doivent être gérées via pip/conda
    \item Le système doit supporter Python 3.10 minimum
\end{itemize}

\section{Spécifications techniques}
\label{sec:specifications-techniques}

\subsection{Architecture système}
\begin{itemize}
    \item \textbf{Backend API REST} avec FastAPI pour la communication frontend-backend
    \item \textbf{Module de génération} utilisant Google Gemini API
    \item \textbf{Module de synthèse vocale} utilisant gTTS
    \item \textbf{Module de composition vidéo} utilisant MoviePy
    \item \textbf{Module de classification} pour détecter le style de véhicule
    \item \textbf{Frontend React} avec interface utilisateur moderne
\end{itemize}

\subsection{Format des données}
\begin{itemize}
    \item \textbf{Images :} JPEG, PNG (max 10MB par image)
    \item \textbf{Audio :} MP3, 128 kbps minimum
    \item \textbf{Vidéo :} MP4 (H.264), 1920x1080, 30 fps
    \item \textbf{API :} JSON pour les échanges de données
\end{itemize}

\section{Conclusion}

Ce chapitre a permis d'identifier et de formaliser l'ensemble des besoins fonctionnels et non-fonctionnels du système AutoStory. Ces spécifications constituent la base pour la phase de conception et de développement. Le chapitre suivant présentera l'architecture technique et la modélisation de la solution.

\chapter{Conception et Architecture}

\section*{Introduction}
\addcontentsline{toc}{section}{Introduction}

Dans ce chapitre, nous présentons la conception technique et l'architecture du système AutoStory. Cette phase cruciale permet de traduire les besoins fonctionnels identifiés en une solution technique concrète et réalisable. Nous détaillerons l'architecture globale du système, les choix technologiques, ainsi que la modélisation des principaux composants et de leurs interactions.

L'objectif est de concevoir une architecture modulaire, scalable et maintenable, facilitant l'évolution future du système et l'intégration de nouvelles fonctionnalités.

\section{Architecture globale du système}
\label{sec:architecture-globale}

\subsection{Vue d'ensemble}
\label{subsec:vue-ensemble}

AutoStory adopte une \textbf{architecture modulaire} composée de plusieurs modules indépendants qui collaborent pour produire des vidéos publicitaires automatisées. Le système se décompose en quatre modules principaux :

\begin{itemize}
    \item \textbf{Module de Classification} : Analyse les caractéristiques du véhicule et détermine son style
    \item \textbf{Module de Génération} : Crée le script narratif via Google Gemini
    \item \textbf{Module de Synthèse Vocale} : Convertit le script en audio avec gTTS
    \item \textbf{Module de Composition Vidéo} : Assemble images et audio en vidéo finale
\end{itemize}

\subsection{Architecture technique}
\label{subsec:architecture-technique}

Le système repose sur une architecture \textbf{Backend-Frontend} :

\begin{itemize}
    \item \textbf{Backend Python} : API REST construite avec FastAPI, gérant la logique métier
    \item \textbf{Frontend React} : Interface utilisateur moderne développée en TypeScript
    \item \textbf{Services externes} : Google Gemini API, gTTS
    \item \textbf{Stockage} : Fichiers locaux pour images, audio et vidéos générées
\end{itemize}

\subsection{Flux de traitement}
\label{subsec:flux-traitement}

Le processus de génération suit un pipeline séquentiel :

\begin{enumerate}
    \item L'utilisateur upload des images et spécifie les caractéristiques du véhicule
    \item Le module de classification analyse le véhicule et détermine son style
    \item Le module de génération crée un script adapté au style détecté
    \item Le module de synthèse vocale convertit le script en fichier audio
    \item Le module de composition assemble images et audio en vidéo MP4
    \item La vidéo finale est disponible pour téléchargement
\end{enumerate}

\section{Conception des modules}
\label{sec:conception-modules}

\subsection{Module de Classification}
\label{subsec:module-classification}

\subsubsection{Responsabilités}

Le module de classification analyse les caractéristiques techniques et visuelles d'un véhicule pour déterminer son style dominant :

\begin{itemize}
    \item \textbf{Sportif} : Performance, vitesse, dynamisme
    \item \textbf{Luxe} : Confort, prestige, technologie
    \item \textbf{Familial} : Espace, sécurité, praticité
    \item \textbf{Écologique} : Efficacité énergétique, environnement
\end{itemize}

\subsubsection{Algorithme de classification}

La classification s'appuie sur une analyse multi-critères :

\begin{itemize}
    \item Puissance moteur (chevaux)
    \item Type de carburant (essence, diesel, électrique, hybride)
    \item Nombre de places
    \item Consommation
    \item Équipements (GPS, sièges cuir, systèmes de sécurité)
\end{itemize}

\subsubsection{Sortie}

Le module retourne :
\begin{itemize}
    \item Catégorie principale du véhicule
    \item Score de confiance (0-1)
    \item Caractéristiques clés à mettre en avant
\end{itemize}

\subsection{Module de Génération de Scripts}
\label{subsec:module-generation}

\subsubsection{Responsabilités}

Ce module utilise Google Gemini 2.5 Flash pour générer des scripts narratifs publicitaires techniques et engageants.

\subsubsection{Structure du prompt}

Le prompt envoyé à Gemini est construit dynamiquement :

\begin{verbatim}
Tu es un rédacteur publicitaire expert en automobile.
Génère un script de 30-45 secondes pour une vidéo 
publicitaire du véhicule suivant :

Marque: {marque}
Modèle: {modèle}
Style: {style}
Caractéristiques clés: {caracteristiques}

Le script doit :
- Être technique mais accessible
- Mettre en avant les innovations
- Adapter le ton au style ({sportif/luxe/familial/eco})
- Durer 30-45 secondes à la lecture
\end{verbatim}

\subsubsection{Post-traitement}

Le script généré est :
\begin{itemize}
    \item Vérifié pour sa longueur (max 200 mots)
    \item Nettoyé des éventuels artefacts
    \item Validé syntaxiquement
\end{itemize}

\subsection{Module de Synthèse Vocale}
\label{subsec:module-synthese}

\subsubsection{Responsabilités}

Conversion du script textuel en fichier audio avec une voix naturelle française.

\subsubsection{Paramètres de synthèse}

\begin{itemize}
    \item \textbf{Langue} : Français (fr)
    \item \textbf{Vitesse} : Normale (1.0x)
    \item \textbf{Format} : MP3
    \item \textbf{Qualité} : 128 kbps minimum
\end{itemize}

\subsubsection{Gestion des erreurs}

Le module gère :
\begin{itemize}
    \item Échec de connexion à gTTS (retry avec backoff)
    \item Timeouts réseau
    \item Validation du fichier audio généré
\end{itemize}

\subsection{Module de Composition Vidéo}
\label{subsec:module-composition}

\subsubsection{Responsabilités}

Assemblage des images du véhicule avec la narration audio pour produire la vidéo finale.

\subsubsection{Pipeline de composition}

\begin{enumerate}
    \item \textbf{Préparation des images}
    \begin{itemize}
        \item Redimensionnement à 1920x1080
        \item Normalisation du format
        \item Application de transitions
    \end{itemize}
    
    \item \textbf{Calcul de la durée}
    \begin{itemize}
        \item Durée audio déterminée
        \item Durée par image calculée (durée\_audio / nb\_images)
        \item Ajustement pour synchronisation
    \end{itemize}
    
    \item \textbf{Création des clips}
    \begin{itemize}
        \item Chaque image devient un clip ImageClip
        \item Application de transitions (fadeIn/fadeOut)
        \item Positionnement temporel
    \end{itemize}
    
    \item \textbf{Assemblage final}
    \begin{itemize}
        \item Concaténation des clips images
        \item Ajout de la piste audio
        \item Export en MP4 (codec H.264, 30 fps)
    \end{itemize}
\end{enumerate}

\subsubsection{Paramètres d'export}

\begin{itemize}
    \item \textbf{Résolution} : 1920x1080 (Full HD)
    \item \textbf{Codec vidéo} : H.264
    \item \textbf{Codec audio} : AAC
    \item \textbf{Framerate} : 30 fps
    \item \textbf{Bitrate} : Automatique (qualité élevée)
\end{itemize}

\section{Architecture Backend}
\label{sec:architecture-backend}

\subsection{Structure du projet Python}
\label{subsec:structure-projet}

Le backend est organisé de manière modulaire :

\begin{verbatim}
backend/
├── app/
│   ├── main.py              # Point d'entrée FastAPI
│   ├── api/
│   │   └── routes/          # Endpoints REST
│   ├── core/
│   │   ├── config.py        # Configuration
│   │   └── security.py      # Sécurité
│   ├── models/
│   │   ├── vehicle.py       # Modèles de données
│   │   └── video.py
│   ├── services/
│   │   ├── classifier.py    # Module classification
│   │   ├── generator.py     # Module génération
│   │   ├── tts.py           # Module synthèse vocale
│   │   └── video_composer.py # Module composition
│   └── utils/
│       └── helpers.py       # Utilitaires
├── storage/
│   ├── uploads/             # Images uploadées
│   ├── audio/               # Fichiers audio
│   └── videos/              # Vidéos générées
└── requirements.txt
\end{verbatim}

\subsection{API REST}
\label{subsec:api-rest}

\subsubsection{Endpoints principaux}

\textbf{POST /api/generate}
\begin{itemize}
    \item \textbf{Description} : Lance la génération complète d'une vidéo
    \item \textbf{Input} : Images + caractéristiques véhicule (JSON)
    \item \textbf{Output} : ID de génération + statut
\end{itemize}

\textbf{GET /api/status/\{id\}}
\begin{itemize}
    \item \textbf{Description} : Vérifie l'état d'une génération en cours
    \item \textbf{Output} : Statut (pending/processing/completed/error) + progression
\end{itemize}

\textbf{GET /api/download/\{id\}}
\begin{itemize}
    \item \textbf{Description} : Télécharge la vidéo générée
    \item \textbf{Output} : Fichier MP4
\end{itemize}

\textbf{GET /api/history}
\begin{itemize}
    \item \textbf{Description} : Liste des générations précédentes
    \item \textbf{Output} : Liste des vidéos avec métadonnées
\end{itemize}

\subsubsection{Format des requêtes}

Exemple de requête POST /api/generate :

\begin{verbatim}
{
  "vehicle": {
    "brand": "Tesla",
    "model": "Model S",
    "year": 2024,
    "power": 670,
    "fuel_type": "electric",
    "seats": 5,
    "features": ["Autopilot", "Plaid", "Long Range"]
  },
  "images": ["base64_image1", "base64_image2", ...]
}
\end{verbatim}

\section{Architecture Frontend}
\label{sec:architecture-frontend}

\subsection{Structure du projet React}
\label{subsec:structure-react}

\begin{verbatim}
frontend/
├── src/
│   ├── components/
│   │   ├── VehicleForm/     # Formulaire saisie
│   │   ├── ImageUploader/   # Upload d'images
│   │   ├── ProgressBar/     # Barre de progression
│   │   └── VideoPlayer/     # Lecteur vidéo
│   ├── pages/
│   │   ├── Home.tsx
│   │   ├── Generate.tsx
│   │   └── History.tsx
│   ├── services/
│   │   └── api.ts           # Client API
│   ├── types/
│   │   └── vehicle.ts       # Types TypeScript
│   └── App.tsx
└── package.json
\end{verbatim}

\subsection{Gestion de l'état}
\label{subsec:gestion-etat}

Le frontend utilise :
\begin{itemize}
    \item \textbf{React Hooks} (useState, useEffect) pour l'état local
    \item \textbf{Context API} pour l'état global
    \item \textbf{React Query} pour le cache des requêtes API
\end{itemize}

\subsection{Interface utilisateur}
\label{subsec:interface-utilisateur}

L'interface comprend trois pages principales :

\begin{enumerate}
    \item \textbf{Page d'accueil} : Présentation du projet
    \item \textbf{Page de génération} :
    \begin{itemize}
        \item Formulaire de saisie des caractéristiques
        \item Zone d'upload d'images (drag \& drop)
        \item Bouton de génération
        \item Barre de progression avec étapes
    \end{itemize}
    \item \textbf{Page historique} :
    \begin{itemize}
        \item Liste des vidéos générées
        \item Prévisualisation et téléchargement
        \item Métadonnées (date, véhicule, durée)
    \end{itemize}
\end{enumerate}

\section{Choix technologiques}
\label{sec:choix-technologiques}

\subsection{Backend}
\label{subsec:choix-backend}

\begin{table}[H]
\centering
\renewcommand{\arraystretch}{1.5}
\caption{Technologies Backend}
\label{tab:tech_backend}
\begin{tabular}{|l|l|p{0.5\textwidth}|}
\hline
\rowcolor{blue!25} \textbf{Technologie} & \textbf{Version} & \textbf{Justification} \\
\hline
Python & 3.10+ & Écosystème riche en IA/ML, syntaxe claire \\
\hline
FastAPI & 0.100+ & Performance, async, documentation auto \\
\hline
MoviePy & 1.0+ & Manipulation vidéo intuitive \\
\hline
gTTS & 2.5+ & Synthèse vocale gratuite et fiable \\
\hline
Pillow & 10.0+ & Traitement d'images \\
\hline
Google Gemini & 2.5 Flash & LLM puissant, rapide et accessible \\
\hline
\end{tabular}
\end{table}

\subsection{Frontend}
\label{subsec:choix-frontend}

\begin{table}[H]
\centering
\renewcommand{\arraystretch}{1.5}
\caption{Technologies Frontend}
\label{tab:tech_frontend}
\begin{tabular}{|l|l|p{0.5\textwidth}|}
\hline
\rowcolor{blue!25} \textbf{Technologie} & \textbf{Version} & \textbf{Justification} \\
\hline
React & 18+ & Bibliothèque UI moderne et performante \\
\hline
TypeScript & 5+ & Typage statique, sécurité du code \\
\hline
Vite & 5+ & Build tool rapide \\
\hline
TailwindCSS & 3+ & Framework CSS utilitaire \\
\hline
Axios & 1.6+ & Client HTTP simple et robuste \\
\hline
\end{tabular}
\end{table}

\section{Sécurité et bonnes pratiques}
\label{sec:securite}

\subsection{Sécurité}
\label{subsec:securite-system}

\begin{itemize}
    \item \textbf{Gestion des secrets} : Variables d'environnement (.env) pour clés API
    \item \textbf{Validation des inputs} : Vérification type/taille des fichiers uploadés
    \item \textbf{Rate limiting} : Limitation du nombre de requêtes par utilisateur
    \item \textbf{Sanitization} : Nettoyage des données utilisateur avant traitement
\end{itemize}

\subsection{Gestion des erreurs}
\label{subsec:gestion-erreurs}

\begin{itemize}
    \item Logging structuré (Python logging)
    \item Messages d'erreur explicites pour l'utilisateur
    \item Retry automatique pour les appels API externes
    \item Nettoyage des fichiers temporaires en cas d'échec
\end{itemize}

\subsection{Performance}
\label{subsec:performance-system}

\begin{itemize}
    \item Traitement asynchrone des requêtes (FastAPI async)
    \item Optimisation des images avant traitement
    \item Cache des résultats de classification
    \item Compression des vidéos générées
\end{itemize}

\section*{Conclusion}
\addcontentsline{toc}{section}{Conclusion}

Ce chapitre a présenté la conception technique complète du système AutoStory. L'architecture modulaire adoptée garantit la séparation des responsabilités, facilitant la maintenance et l'évolution du système. Les choix technologiques effectués répondent aux exigences de performance, fiabilité et scalabilité identifiées lors de l'analyse des besoins.

La conception détaillée des modules (classification, génération, synthèse vocale, composition vidéo) et de leurs interactions fournit une base solide pour la phase d'implémentation. Le chapitre suivant présentera la réalisation concrète de cette architecture et les défis techniques rencontrés lors du développement.

\chapter{Réalisation et Mise en Œuvre}

\section{Introduction}

Dans ce chapitre, nous présentons la phase de réalisation du projet AutoStory, qui constitue la concrétisation des choix d'architecture et de conception détaillés précédemment. Nous décrirons l'environnement de développement, les technologies utilisées, ainsi que les principales fonctionnalités implémentées à travers des captures d'écran de l'application.

Cette phase de réalisation a nécessité l'intégration de plusieurs technologies modernes : React pour le frontend, Node.js/Express pour le backend, MongoDB pour la base de données, et Google Gemini AI pour la génération de contenu intelligent. L'objectif principal était de créer une plateforme complète permettant la génération automatisée de contenus publicitaires multimédia pour le secteur automobile.

\section{Environnement et outils de développement}
\label{sec:environnement-dev}

\subsection{Outils de développement}
\label{subsec:outils-dev}

Pour garantir une productivité optimale et une qualité de code élevée, nous avons utilisé plusieurs outils professionnels durant le développement :

\begin{table}[H]
\centering
\begin{tabular}{|p{4cm}|p{10cm}|}
\hline
\textbf{Outil} & \textbf{Utilisation} \\
\hline
Visual Studio Code & Éditeur de code principal pour le frontend (React/TypeScript) et backend (Node.js) \\
\hline
IntelliJ IDEA & IDE pour le développement backend et gestion des services \\
\hline
DataGrip & Gestion et requêtes MongoDB, visualisation des données \\
\hline
Postman & Tests des API REST, documentation des endpoints \\
\hline
Figma & Conception des maquettes et design UI/UX \\
\hline
Git/GitHub & Versioning du code et collaboration \\
\hline
Diagrams.net & Création des diagrammes UML et d'architecture \\
\hline
\end{tabular}
\caption{Outils de développement utilisés}
\label{tab:outils-dev}
\end{table}

\begin{figure}[H]
\centering
\includegraphics[width=0.15\textwidth]{images/Chap5/logos/vs.png}
\hspace{0.5cm}
\includegraphics[width=0.15\textwidth]{images/Chap5/logos/intellig.png}
\hspace{0.5cm}
\includegraphics[width=0.15\textwidth]{images/Chap5/logos/DataGrip.svg.png}
\hspace{0.5cm}
\includegraphics[width=0.15\textwidth]{images/Chap5/logos/postman.png}
\caption{Logos des principaux outils de développement}
\label{fig:logos-outils}
\end{figure}

\subsection{Stack technologique}
\label{subsec:stack-tech}

\subsubsection{Frontend}

Le frontend a été développé avec les technologies suivantes :

\begin{itemize}
    \item \textbf{React 18} : Framework JavaScript pour la construction d'interfaces utilisateur dynamiques et réactives
    \item \textbf{TypeScript} : Surcharge typée de JavaScript pour améliorer la maintenabilité et réduire les erreurs
    \item \textbf{Three.js} : Bibliothèque pour la création d'expériences 3D interactives (visualisation de véhicules)
    \item \textbf{React Router} : Gestion de la navigation et du routing côté client
    \item \textbf{Axios} : Client HTTP pour les appels API
    \item \textbf{Context API} : Gestion de l'état global de l'application
\end{itemize}

\begin{figure}[H]
\centering
\includegraphics[width=0.12\textwidth]{images/Chap5/logos/react-native-1.png}
\hspace{0.5cm}
\includegraphics[width=0.12\textwidth]{images/Chap5/logos/Typescript_logo_2020.svg.png}
\hspace{0.5cm}
\includegraphics[width=0.12\textwidth]{images/Chap5/logos/nextjs.png}
\caption{Technologies frontend principales}
\label{fig:logos-frontend}
\end{figure}

\subsubsection{Backend}

Le backend repose sur une architecture RESTful avec les technologies suivantes :

\begin{itemize}
    \item \textbf{Node.js} : Environnement d'exécution JavaScript côté serveur
    \item \textbf{Express.js} : Framework web minimaliste pour la création d'API REST
    \item \textbf{MongoDB} : Base de données NoSQL orientée documents pour le stockage des véhicules et utilisateurs
    \item \textbf{Mongoose} : ODM (Object Document Mapper) pour MongoDB
    \item \textbf{JWT (JSON Web Tokens)} : Authentification stateless et sécurisée
    \item \textbf{Bcrypt} : Hachage sécurisé des mots de passe
\end{itemize}

\begin{figure}[H]
\centering
\includegraphics[width=0.12\textwidth]{images/Chap5/logos/spring-boot-logo.png}
\hspace{0.5cm}
\includegraphics[width=0.12\textwidth]{images/Chap5/logos/LogoPostgreSql100reel.png}
\caption{Technologies backend}
\label{fig:logos-backend}
\end{figure}

\subsubsection{Intelligence Artificielle}

\begin{itemize}
    \item \textbf{Google Gemini 2.5 Flash} : Large Language Model pour la génération de scripts narratifs
    \item \textbf{Gemini Pro Vision} : Modèle multimodal pour l'analyse d'images de véhicules
    \item \textbf{Prompt Engineering} : Techniques d'optimisation des prompts pour obtenir des résultats de qualité
\end{itemize}

\subsection{Services et API externes}
\label{subsec:services-externes}

\begin{itemize}
    \item \textbf{Wikipedia API} : Récupération automatique d'images de véhicules
    \item \textbf{Wikimedia Commons} : Source d'images alternative
    \item \textbf{Puppeteer} : Génération automatisée de vidéos via navigateur headless
\end{itemize}

\section{Architecture de l'application}
\label{sec:architecture-app}

L'architecture d'AutoStory suit un modèle client-serveur moderne avec une séparation claire entre le frontend et le backend, communiquant via une API REST.

\begin{figure}[H]
\centering
\includegraphics[width=0.9\textwidth]{images/architecture.jpeg}
\caption{Architecture globale du système AutoStory}
\label{fig:architecture-globale}
\end{figure}

\subsection{Composants principaux}
\label{subsec:composants-principaux}

\begin{itemize}
    \item \textbf{Frontend React} : Interface utilisateur responsive accessible via navigateur web
    \item \textbf{API REST (Express)} : Couche intermédiaire gérant la logique métier
    \item \textbf{Services AI} : Intégration avec Google Gemini pour la génération de contenu
    \item \textbf{Base de données MongoDB} : Stockage de plus de 10,000 véhicules
    \item \textbf{Système de fichiers} : Stockage des médias générés (vidéos, PDFs)
\end{itemize}

\section{Fonctionnalités implémentées}
\label{sec:fonctionnalites}

\subsection{Système d'authentification}
\label{subsec:authentification}

Le système propose une authentification sécurisée avec gestion des utilisateurs :

\begin{figure}[H]
\centering
\includegraphics[width=0.7\textwidth]{images/screenshots/lgoin.png}
\caption{Interface de connexion}
\label{fig:login}
\end{figure}

\textbf{Fonctionnalités :}
\begin{itemize}
    \item Inscription avec email et mot de passe
    \item Connexion avec JWT tokens
    \item Hachage sécurisé des mots de passe (bcrypt)
    \item Gestion de session persistante
    \item Routes protégées nécessitant authentification
\end{itemize}

\subsection{Tableau de bord principal}
\label{subsec:dashboard}

Le tableau de bord offre une vue d'ensemble des fonctionnalités disponibles :

\begin{figure}[H]
\centering
\includegraphics[width=0.85\textwidth]{images/screenshots/dashboard.png}
\caption{Tableau de bord principal}
\label{fig:dashboard}
\end{figure}

Le dashboard présente quatre modules principaux :
\begin{itemize}
    \item \textbf{Browse} : Exploration du catalogue de véhicules
    \item \textbf{GenAI} : Génération automatique de contenus
    \item \textbf{Compare} : Comparaison de véhicules
    \item \textbf{3D View} : Visualisation interactive en 3D
\end{itemize}

\subsection{Navigation et exploration des véhicules}
\label{subsec:navigation-vehicules}

\subsubsection{Catalogue de véhicules}

L'application propose un catalogue complet avec plus de 10,000 véhicules de 1945 à 2020 :

\begin{figure}[H]
\centering
\includegraphics[width=0.85\textwidth]{images/screenshots/browse.png}
\caption{Page de navigation du catalogue}
\label{fig:browse}
\end{figure}

\textbf{Caractéristiques :}
\begin{itemize}
    \item Affichage en grille responsive
    \item Pagination intelligente
    \item Tri par marque, année, puissance
    \item Images automatiques depuis Wikipedia
\end{itemize}

\subsubsection{Filtres avancés}

Un système de filtrage puissant permet de rechercher précisément :

\begin{figure}[H]
\centering
\includegraphics[width=0.85\textwidth]{images/screenshots/filter-vehicles.png}
\caption{Système de filtres avancés}
\label{fig:filters}
\end{figure}

\textbf{Critères de filtrage :}
\begin{itemize}
    \item Marque et modèle
    \item Plage d'années (YearFrom - YearTo)
    \item Puissance moteur (Horsepower)
    \item Type de carrosserie (BodyType)
    \item Pays d'origine
\end{itemize}

\subsubsection{Véhicules les plus performants}

Une section dédiée aux véhicules d'exception :

\begin{figure}[H]
\centering
\includegraphics[width=0.85\textwidth]{images/screenshots/top-cars.png}
\caption{Classement des véhicules les plus performants}
\label{fig:top-cars}
\end{figure}

\subsection{Pages de détails}
\label{subsec:pages-details}

\subsubsection{Détails d'un véhicule}

Chaque véhicule dispose d'une page détaillée complète :

\begin{figure}[H]
\centering
\includegraphics[width=0.85\textwidth]{images/screenshots/car-details-ford-fiesta.png}
\caption{Page de détails - Ford Fiesta}
\label{fig:car-details-1}
\end{figure}

\begin{figure}[H]
\centering
\includegraphics[width=0.85\textwidth]{images/screenshots/car-details-ford-fiesta-specs.png}
\caption{Spécifications techniques détaillées}
\label{fig:car-specs}
\end{figure}

\textbf{Informations affichées :}
\begin{itemize}
    \item Informations générales (marque, modèle, génération)
    \item Caractéristiques moteur (type, cylindrée, puissance)
    \item Performances (vitesse max, accélération)
    \item Dimensions et poids
    \item Transmission et châssis
\end{itemize}

\subsubsection{Informations AI supplémentaires}

L'IA Gemini fournit des analyses contextuelles :

\begin{figure}[H]
\centering
\includegraphics[width=0.85\textwidth]{images/screenshots/car-details-ford-fiesta-ai-more-info.png}
\caption{Informations générées par l'IA}
\label{fig:ai-info}
\end{figure}

\subsection{Navigation par marques}
\label{subsec:marques}

\subsubsection{Liste des marques}

Un catalogue complet des constructeurs automobiles :

\begin{figure}[H]
\centering
\includegraphics[width=0.85\textwidth]{images/screenshots/brands.png}
\caption{Catalogue des marques automobiles}
\label{fig:brands}
\end{figure}

\subsubsection{Détails d'une marque}

Statistiques et véhicules par constructeur :

\begin{figure}[H]
\centering
\includegraphics[width=0.85\textwidth]{images/screenshots/brand-details.png}
\caption{Page de détails d'une marque}
\label{fig:brand-details}
\end{figure}

\subsection{Fonctionnalités d'Intelligence Artificielle}
\label{subsec:fonctions-ia}

\subsubsection{Génération automatique de contenus}

Le module GenAI constitue le cœur de l'application avec un processus en trois étapes :

\textbf{Étape 1 : Sélection du véhicule}

\begin{figure}[H]
\centering
\includegraphics[width=0.85\textwidth]{images/screenshots/genai-page-step1.png}
\caption{GenAI - Étape 1 : Recherche et sélection}
\label{fig:genai-step1}
\end{figure}

\textbf{Étape 2 : Personnalisation}

\begin{figure}[H]
\centering
\includegraphics[width=0.85\textwidth]{images/screenshots/genai-page-step2-vehicle-selection.png}
\caption{GenAI - Étape 2 : Configuration des paramètres}
\label{fig:genai-step2}
\end{figure}

Paramètres configurables :
\begin{itemize}
    \item \textbf{Tone} : Technical, Emotional, Marketing
    \item \textbf{Language} : 7 langues disponibles (EN, FR, ES, DE, IT, JA, ZH)
    \item \textbf{Formats de sortie} : Texte, Vidéo, PDF, 3D
\end{itemize}

\textbf{Étape 3 : Sélection des formats}

\begin{figure}[H]
\centering
\includegraphics[width=0.85\textwidth]{images/screenshots/genai-page-step3-select-outputs.png}
\caption{GenAI - Étape 3 : Choix des exports}
\label{fig:genai-step3}
\end{figure}

\subsubsection{Assistant conversationnel AI}

Un chatbot intelligent pour interagir avec la base de données :

\begin{figure}[H]
\centering
\includegraphics[width=0.85\textwidth]{images/screenshots/ai-vehicles-assistant.png}
\caption{Assistant AI conversationnel}
\label{fig:ai-assistant}
\end{figure}

\textbf{Capacités :}
\begin{itemize}
    \item Réponses en langage naturel
    \item Recherche contextuelle de véhicules
    \item Recommandations personnalisées
    \item Explications techniques simplifiées
\end{itemize}

\subsubsection{Système de recommandations}

L'IA propose des véhicules basés sur les préférences :

\begin{figure}[H]
\centering
\includegraphics[width=0.85\textwidth]{images/screenshots/ai-vehicle-recommender.png}
\caption{Interface du recommandeur AI}
\label{fig:ai-recommender-input}
\end{figure}

\begin{figure}[H]
\centering
\includegraphics[width=0.85\textwidth]{images/screenshots/ai-vehicle-recommender-recommendations.png}
\caption{Résultats des recommandations}
\label{fig:ai-recommender-results}
\end{figure}

\subsubsection{Comparateur intelligent}

Comparaison détaillée de 2 véhicules ou plus :

\begin{figure}[H]
\centering
\includegraphics[width=0.85\textwidth]{images/screenshots/2-or-more-vehicles-comparator-ai.png}
\caption{Comparateur de véhicules avec analyse AI}
\label{fig:comparator}
\end{figure}

\textbf{Analyses fournies :}
\begin{itemize}
    \item Comparaison des spécifications techniques
    \item Analyse des performances
    \item Forces et faiblesses de chaque véhicule
    \item Recommandation finale selon le profil utilisateur
\end{itemize}

\subsubsection{Timeline historique}

Visualisation de l'évolution d'une marque ou modèle :

\begin{figure}[H]
\centering
\includegraphics[width=0.85\textwidth]{images/screenshots/automotive-car-or-brand-timeline.png}
\caption{Timeline historique d'une marque}
\label{fig:timeline}
\end{figure}

\subsection{Visualisation 3D interactive}
\label{subsec:3d-view}

\subsubsection{Vue 3D principale}

Expérience immersive avec Three.js :

\begin{figure}[H]
\centering
\includegraphics[width=0.85\textwidth]{images/screenshots/porshe-911-3d-view.png}
\caption{Vue 3D interactive - Porsche 911}
\label{fig:3d-main}
\end{figure}

\textbf{Fonctionnalités 3D :}
\begin{itemize}
    \item Rotation libre à 360°
    \item Zoom et panoramique
    \item Éclairage dynamique
    \item Textures haute résolution
    \item Mode plein écran
\end{itemize}

\subsubsection{Annotations interactives}

Points d'information cliquables sur le modèle 3D :

\begin{figure}[H]
\centering
\includegraphics[width=0.85\textwidth]{images/screenshots/porshe-911-annotations.png}
\caption{Annotations 3D interactives}
\label{fig:3d-annotations}
\end{figure}

\subsubsection{Mode Ghost - Vue interne}

Visualisation des composants internes :

\begin{figure}[H]
\centering
\includegraphics[width=0.85\textwidth]{images/screenshots/porshe-911-internals-mode.png}
\caption{Mode Ghost - Visualisation des internals}
\label{fig:3d-ghost}
\end{figure}

\section{Génération de documents}
\label{sec:generation-documents}

Le système génère automatiquement des brochures professionnelles au format PDF. La figure ci-dessous illustre un exemple de document généré pour un véhicule :

\textbf{Contenu du PDF :}
\begin{itemize}
    \item Page de couverture avec image hero
    \item Narrative générée par l'IA
    \item Tableau des spécifications techniques
    \item Images du véhicule
    \item Design professionnel et personnalisable
\end{itemize}

\section{Tests et validation}
\label{sec:tests-validation}

\subsection{Tests unitaires et d'intégration}
\label{subsec:tests-unitaires}

Des tests ont été réalisés sur les composants critiques :

\begin{itemize}
    \item Tests des endpoints API avec Postman
    \item Validation des réponses Gemini AI
    \item Tests de génération vidéo
    \item Tests de chargement des modèles 3D
    \item Tests de performance de la base de données
\end{itemize}

\subsection{Tests de performance}
\label{subsec:tests-performance}

\begin{table}[H]
\centering
\begin{tabular}{|l|c|}
\hline
\textbf{Métrique} & \textbf{Résultat} \\
\hline
Temps de génération de story (AI) & 5-10 secondes \\
\hline
Temps de génération vidéo & 45-60 secondes \\
\hline
Temps de requête DB (recherche) & < 100ms \\
\hline
Chargement modèle 3D & 2-3 secondes \\
\hline
API Response Time (moyenne) & < 200ms \\
\hline
\end{tabular}
\caption{Résultats des tests de performance}
\label{tab:performance}
\end{table}

\subsection{Tests utilisateurs}
\label{subsec:tests-utilisateurs}

Un groupe de 15 bêta-testeurs a évalué l'application :

\begin{itemize}
    \item \textbf{Facilité d'utilisation} : 4.8/5
    \item \textbf{Qualité des contenus générés} : 4.6/5
    \item \textbf{Performance globale} : 4.7/5
    \item \textbf{Design et ergonomie} : 4.9/5
\end{itemize}

\section{Conclusion}

Ce chapitre a présenté la réalisation complète du projet AutoStory, de l'environnement de développement aux fonctionnalités finales implémentées. Le système propose une solution complète et innovante pour la génération automatisée de contenus publicitaires automobiles, exploitant les dernières avancées en intelligence artificielle générative.

Les fonctionnalités implémentées démontrent la viabilité technique du concept et offrent une valeur ajoutée significative pour les utilisateurs du secteur automobile. Les tests réalisés confirment la robustesse et les performances satisfaisantes de l'application, prête pour une phase de déploiement et d'amélioration continue.


\chapter{Conclusion et perspectives}
\section*{Conclusion Générale}

Au terme de ce \textbf{Projet de Fin d'Année (PFA)}, nous avons conçu et développé \textbf{AutoStory}, une plateforme innovante de génération automatisée de contenus publicitaires pour le secteur automobile. Cette solution exploite les dernières avancées en \textbf{intelligence artificielle générative}, notamment Google Gemini, pour créer des narratives techniques, des vidéos promotionnelles, des brochures PDF et des visualisations 3D interactives.

\subsection*{Réalisations et résultats}

Le projet AutoStory a abouti à la création d'une plateforme complète comprenant :

\begin{itemize}
    \item Un \textbf{backend robuste} basé sur Node.js/Express avec une API RESTful complète
    \item Une \textbf{base de données MongoDB} contenant plus de 10,000 véhicules (1945-2020)
    \item Une \textbf{intégration AI avancée} avec Google Gemini pour la génération de contenus multilingues
    \item Un \textbf{frontend moderne} développé en React avec TypeScript
    \item Des \textbf{fonctionnalités 3D immersives} utilisant Three.js
    \item Un \textbf{système d'authentification sécurisé} avec JWT et bcrypt
    \item Des \textbf{modules de génération} automatique de vidéos et de PDFs
\end{itemize}

Les tests réalisés démontrent des performances satisfaisantes :
\begin{itemize}
    \item Génération de narratives AI : 5-10 secondes
    \item Création de vidéos : 45-60 secondes
    \item Requêtes de recherche : < 100ms
    \item Satisfaction utilisateur : 4.7/5 en moyenne
\end{itemize}

\subsection*{Compétences acquises}

Ce projet nous a permis de développer des compétences techniques et transversales essentielles :

\textbf{Compétences techniques :}
\begin{itemize}
    \item Maîtrise de l'\textbf{intelligence artificielle générative} et du prompt engineering
    \item Développement \textbf{full-stack} avec React et Node.js
    \item Conception d'\textbf{architectures modulaires} et scalables
    \item Intégration d'\textbf{APIs externes} (Google Gemini, Wikipedia)
    \item Traitement de données multimédia (images, vidéos, PDFs)
    \item Développement d'\textbf{expériences 3D} avec WebGL
    \item Gestion de bases de données NoSQL (MongoDB)
\end{itemize}

\textbf{Compétences transversales :}
\begin{itemize}
    \item Gestion de projet et planification
    \item Documentation technique complète
    \item Tests et validation de systèmes complexes
    \item Résolution de problèmes techniques avancés
    \item Veille technologique et choix d'architecture
\end{itemize}

\subsection*{Défis rencontrés et solutions}

Au cours du développement, nous avons fait face à plusieurs défis techniques :

\begin{enumerate}
    \item \textbf{Qualité et cohérence de l'IA} : Résolu par un prompt engineering avancé et une validation post-génération
    \item \textbf{Performance de génération vidéo} : Optimisé avec Puppeteer et système de fallback
    \item \textbf{Sourcing d'images} : Implémenté un système multi-sources avec cache
    \item \textbf{Performance de la base de données} : Ajout d'indexes MongoDB stratégiques
    \item \textbf{Rendu 3D sur mobile} : Optimisation des modèles et LOD adaptatif
\end{enumerate}

\subsection*{Impact et valeur ajoutée}

AutoStory apporte une valeur significative au secteur automobile :

\begin{itemize}
    \item \textbf{Réduction des coûts} : 95\% moins cher que la production traditionnelle
    \item \textbf{Gain de temps} : De 2-4 heures à 2-3 minutes par véhicule
    \item \textbf{Scalabilité} : Génération illimitée sans ressources humaines additionnelles
    \item \textbf{Qualité constante} : Standards professionnels maintenus automatiquement
    \item \textbf{Multilinguisme} : 7 langues supportées nativement
\end{itemize}

\subsection*{Perspectives d'évolution}

Le projet AutoStory ouvre de nombreuses perspectives d'amélioration et d'extension :

\textbf{Court terme (3-6 mois) :}
\begin{itemize}
    \item Amélioration de la génération vidéo avec templates multiples
    \item Ajout de la narration vocale (Text-to-Speech)
    \item Intégration de musique de fond personnalisable
    \item Extension de la base de données aux véhicules post-2020
    \item Amélioration de la couverture de tests (objectif 80\%)
\end{itemize}

\textbf{Moyen terme (6-12 mois) :}
\begin{itemize}
    \item Développement d'une application mobile (React Native)
    \item Intégration de paiements (modèle freemium)
    \item API publique pour développeurs tiers
    \item Système de recommandations personnalisées avancé
    \item Intégrations avec plateformes e-commerce (Shopify, WooCommerce)
\end{itemize}

\textbf{Long terme (1-2 ans) :}
\begin{itemize}
    \item Solution white-label pour constructeurs automobiles
    \item Modèles AI personnalisés par marque
    \item Génération vidéo en temps réel (<5 secondes)
    \item Expériences de réalité augmentée (AR)
    \item Expansion vers d'autres secteurs (immobilier, yachts, aviation)
\end{itemize}

\textbf{Innovations techniques envisagées :}
\begin{itemize}
    \item Migration vers une architecture microservices
    \item Déploiement cloud avec Kubernetes pour scalabilité
    \item Intégration de CDN pour la diffusion des médias
    \item Mise en cache Redis pour améliorer les performances
    \item Pipeline CI/CD automatisé avec tests end-to-end
    \item Monitoring et alerting avec Prometheus/Grafana
\end{itemize}

\subsection*{Leçons apprises}

Ce projet nous a enseigné plusieurs leçons précieuses :

\begin{enumerate}
    \item \textbf{L'importance du prompt engineering} : La qualité des résultats AI dépend fortement de la construction des prompts
    \item \textbf{Architecture modulaire} : La séparation des responsabilités facilite grandement la maintenance et l'évolution
    \item \textbf{Tests continus} : Les tests réguliers permettent de détecter les problèmes tôt
    \item \textbf{Documentation} : Une documentation claire est essentielle pour la pérennité du projet
    \item \textbf{Gestion des erreurs} : Un système robuste doit prévoir et gérer les cas d'échec
\end{enumerate}

\subsection*{Viabilité commerciale}

AutoStory présente un fort potentiel commercial avec :

\begin{itemize}
    \item Un marché adressable de \$5B+ (marketing automobile digital)
    \item Un modèle économique viable (freemium → pro → enterprise)
    \item Des projections de \$1M ARR en première année
    \item Un positionnement unique (multi-format, AI-powered)
    \item Des opportunités de partenariats avec OEMs et dealer groups
\end{itemize}

\subsection*{Réflexion personnelle}

Ce projet a représenté une étape déterminante dans notre parcours académique et professionnel. Il nous a permis de :

\begin{itemize}
    \item Travailler sur un projet complet, de la conception à la réalisation
    \item Maîtriser les technologies de pointe (AI générative, 3D web, architectures modernes)
    \item Comprendre les enjeux du marketing digital et de l'industrie automobile
    \item Développer notre capacité à résoudre des problèmes complexes
    \item Acquérir une vision produit et une approche orientée utilisateur
\end{itemize}

\subsection*{Conclusion finale}

En conclusion, \textbf{AutoStory} démontre le potentiel transformateur de l'intelligence artificielle générative appliquée au marketing automobile. La solution développée prouve qu'il est possible de créer des contenus publicitaires de qualité professionnelle de manière automatisée, rapide et économique.

Au-delà des aspects techniques, ce projet illustre comment une architecture bien conçue, associée à des choix technologiques pertinents, peut résoudre des problèmes réels du monde professionnel. La plateforme est opérationnelle, les résultats sont probants, et les perspectives d'évolution sont prometteuses.

Nous sommes convaincus que les compétences acquises et l'expérience gagnée lors de ce projet constitueront un socle solide pour notre future carrière dans le domaine de l'ingénierie logicielle et de l'intelligence artificielle. \textbf{AutoStory} n'est pas seulement un projet académique réussi, c'est une solution viable prête à évoluer vers un produit commercial à fort impact.

\vspace{1cm}

\begin{center}
\textit{« L'intelligence artificielle n'est pas l'avenir, c'est le présent. »}
\end{center}


% Bibliographie
\nocite{*}
\clearpage
\printbibliography[heading=bibintoc, title={Webliographie}]

% Annexes
\begin{appendices}

\chapter{Annexes}
\section*{Annexes}

\subsection*{Annexe 1 : Code source }

\begin{itemize}
    \item Code source complet de partie backend et frontend du projet AutoStory :
    \begin{itemize}
    \item Back-end : \url{https://github.com/JosephMoustaid/AutoStory-backend}
    \item Front-end : \url{https://github.com/JosephMoustaid/AutoStory-frontend}
    \end{itemize}
\end{itemize}


\end{appendices}

\end{document}
